% No additional LaTeX package is required for this section.
\subsection{Causal Idle-Time Prediction}
\label{sec:causal-predictor}

The objective of the predictor is not to estimate the most likely idle
duration. Instead, it constructs a \emph{causal lower bound} on the earliest
cycle at which the SA can next be required. At cycle $t$, a certificate of
$H(t)$ cycles asserts
\begin{equation}
  t_{\mathrm{next\mbox{-}SA}} - t \geq H(t),
  \label{eq:causal-certificate}
\end{equation}
using only architectural state visible at or before $t$. The predictor does
not use the subsequently observed idle duration, future DRAM latency, or an
oracle trace. Underestimation only loses a power-saving opportunity;
overestimation could wake the SA late and is therefore treated as a safety
error.

\subsubsection{Visible Next-SA Candidates}

For each of the two resident tiles, the predictor identifies the first
remaining instruction that can use the SA, i.e., the first
\texttt{GEMM\_PRELOAD} or \texttt{GEMM}. The corresponding descriptor stores
the candidate opcode, the number of preceding non-SA instructions, the total
number of remaining tile instructions, and the state of at most two unique
source operands. In the current ONNXim implementation, immutable SA
descriptors are constructed once, after the tile instruction sequence has
been generated and before core admission. A linear pass records every SA
opcode and the number of non-SA instructions since the preceding SA
instruction. During execution, an index selects the current descriptor and
the prefix count decreases as preceding instructions issue. The simulator
then supplies the current descriptor fields to the replayed RTL controller.

The synthesized controller contains the descriptor registers, bound
arithmetic, certificate pipeline, countdown, and clock-control logic. It does
not contain a synthesized hardware implementation of the simulator's
descriptor-generation pass. Therefore, the reported controller area and power do
not include the upstream logic required to generate every descriptor field
from a hardware instruction queue. This is not treated as a zero-cost
operation. A physical implementation would require a compiler, instruction
generator, or hardware producer to provide equivalent fields, but that
producer and its cost are outside the evaluated RTL boundary. We report this
boundary explicitly and do not claim that the synthesized controller
represents the complete metadata-generation overhead.

Let $P_i$ be the number of non-SA instructions before resident tile $i$'s
first SA candidate. Because a core successfully issues at most one resident
instruction per cycle, the candidate cannot issue within the next $P_i$
cycles. If $P_i>0$, the tile-level lower bound is therefore
\begin{equation}
  C_i=P_i.
  \label{eq:prefix-bound-overview}
\end{equation}
If $P_i=0$, the candidate is already at the tile head and its lower bound is
instead derived from the readiness of its required operands. For operand
$j$, the controller combines the core-FIFO and outstanding-response bounds,
denoted $B^{\mathrm{fifo}}_{i,j}$ and
$B^{\mathrm{resp}}_{i,j}$, as
\begin{equation}
  B_{i,j}=\max\!\left(B^{\mathrm{fifo}}_{i,j},
                       B^{\mathrm{resp}}_{i,j}\right).
  \label{eq:operand-combined-bound}
\end{equation}
Both operands must be ready before execution, so the head-candidate bound is
\begin{equation}
  C_i=\max_j B_{i,j}, \qquad P_i=0.
  \label{eq:head-candidate-bound}
\end{equation}
The detailed construction of the prefix, FIFO-ticket, and response-drain
bounds is presented in Section~\ref{sec:causal-bound-components}.

The earliest visible SA use is the minimum over all committed or resident
candidates:
\begin{equation}
  V=\min\!\left(C_0,C_1,C_{\mathrm{pipe}},C_{\mathrm{ex}}\right),
  \label{eq:visible-frontier}
\end{equation}
where $C_{\mathrm{pipe}}$ represents an already committed future SA start and
$C_{\mathrm{ex}}=0$ when the execution-queue head is an SA instruction. A
zero bound conservatively prevents power-state entry.

\subsubsection{Unknown-Future-Work Cap}

A bound derived only from resident tiles must not extend beyond the point at
which an unknown scheduler tile could replace a resident tile. Let $N_i$ be
the number of instructions remaining in resident tile $i$. Since at most one
resident instruction is issued successfully per cycle, tile $i$ cannot retire
before
\begin{equation}
  R_i=\max(1,N_i,C_i).
  \label{eq:retirement-bound}
\end{equation}
If a tile has no remaining SA candidate, the $C_i$ term is omitted and its
retirement bound is $\max(1,N_i)$.
The first resident slot that can become free limits the prediction horizon:
\begin{equation}
  U=\min(R_0,R_1).
  \label{eq:future-tile-cap}
\end{equation}
If fewer than two tiles are resident, an unknown tile may be admitted on the
next cycle and the implementation sets $U=1$. The ordinary causal horizon is
then
\begin{equation}
  H=\begin{cases}
      \min(V,U), & \text{if a visible SA candidate exists},\\
      U,         & \text{otherwise}.
    \end{cases}
  \label{eq:safe-horizon}
\end{equation}
This cap is a conservatism mechanism: it can shorten a valid opportunity, but
it prevents unseen future work from extending an OFF interval.

The implemented controller contains one additional work-drain case. When the
one-cycle unknown-tile cap is the only binding limit and no tile is waiting at
the scheduler, the controller may replace that cap with a global lower bound
on the remaining work that must drain before another tile can be admitted.
If this extension is used, entry is cancelled if scheduler work appears while
the three-stage certificate is being formed.

\subsubsection{Non-Retentive Eligibility}

The target SA loses its resident weight state while OFF. Therefore, a long
idle-time bound alone is insufficient. Entry is considered only when
\begin{enumerate}
  \item the SA is inactive;
  \item the SA compute pipeline and execution queue are empty;
  \item both resident slots are occupied;
  \item the first SA candidate in each resident tile exists and is
        \texttt{GEMM\_PRELOAD}; and
  \item core request-FIFO ticket ordering remains exact.
\end{enumerate}
Requiring \texttt{GEMM\_PRELOAD} ensures that every visible next SA user
reloads the weights before computation. A visible \texttt{GEMM} would depend
on the weights currently retained in the array and therefore disables
non-retentive entry.

The final RTL forms the certificate through three registered stages. Its
electrical break-even parameter is 166 cycles and its wake guard is one
cycle. Consequently, a fresh bound must contain at least
\begin{equation}
  166 + 1 + 3 = 170 \text{ cycles}
  \label{eq:pipelined-entry-threshold}
\end{equation}
before entering the three-stage certificate pipeline. The pipeline subtracts
its three elapsed cycles, and the control domain commits sleep only if at
least $166+1=167$ certified cycles remain. Once committed, the deadline is
frozen: the controller only decrements the saved countdown and never extends
it using newer observations. It wakes immediately if SA activity appears, a
resident descriptor changes, the SA compute pipeline or execution queue
becomes nonempty, ticket ordering becomes inexact, or the countdown reaches
the one-cycle wake guard.

This design separates prediction from observation. Measured future idle
lengths are used only after execution to evaluate coverage and detect safety
violations; they never participate in the certificate used to request the
power-state transition.
